\begin{thebibliography}{99}
\bibitem[1]{ref1} Eric Allender, Anna Gál, and Ian Mertz. Dual VP classes. Comput. Complex., 26(3):583–625, 2017.
\bibitem[2]{ref2} David A. Mix Barrington. Bounded-width polynomial-size branching programs recognize exactly those languages in nc1. J. Comput. Syst. Sci., 38(1):150–164, 1989.
\bibitem[3]{ref3} Michael Ben-Or and Richard Cleve. Computing algebraic formulas using a constant number of registers. SIAM J. Comput., 21(1):54–58, 1992.
\bibitem[4]{ref4} Harry Buhrman, Richard Cleve, Michal Koucký, Bruno Loff, and Florian Speelman. Computing with a full memory: catalytic space. In Symposium on Theory of Computing, STOC 2014, New York, NY, USA, May 31 - June 03, 2014, pages 857–866. ACM, 2014.
\bibitem[5]{ref5} James Cook and Ian Mertz. Catalytic approaches to the tree evaluation problem. In Proccedings of the 52nd Annual ACM SIGACT Symposium on Theory of Computing, STOC 2020, Chicago, IL, USA, June 22-26, 2020, pages 752–760. ACM, 2020.
\bibitem[6]{ref6} James Cook and Ian Mertz. Encodings and the tree evaluation problem. Electron. Colloquium Comput. Complex., page 54, 2021. URL: https://eccc.weizmann.ac.il/report/2021/054.
\bibitem[7]{ref7} Stephen A. Cook, Pierre McKenzie, Dustin Wehr, Mark Braverman, and Rahul Santhanam. Pebbles and branching programs for tree evaluation. ACM Trans. Comput. Theory, 3(2):4:1– 4:43, 2012.
\bibitem[8]{ref8} Vincent Girard, Michal Koucky, and Pierre McKenzie. Nonuniform catalytic space and the direct sum for space. Electronic Colloquium on Computational Complexity (ECCC), 138, 2015.
\bibitem[9]{ref9} Frank Gray. Pulse code communication. https://patents.google.com/patent/US2632058A/ en, 1953. US Patent 2632058A.
\bibitem[10]{ref10} E.I. Nečiporuk. A boolean function. Dokl. Akad. Nauk SSSR, 169(4), 1966.
\bibitem[11]{ref11} Aaron Potechin. A note on amortized branching program complexity, 2017.
\bibitem[12]{ref12} Robert Robere and Jeroen Zuiddam. Amortized circuit complexity, formal complexity measures, and catalytic algorithms. In FOCS, pages 759–769. IEEE, 2021.
\end{thebibliography}
